\begin{abstract}
Agents that revisit evolving data can repeat the same procedural mistake even as the underlying values change: using evidence released after a cutoff, confusing publication time with reporting time, or failing to ground a numerical change in contemporaneous documents. A reusable temporal procedure can persist the missing operation, but a standard targeted-versus-generic comparison can be misleading when the generic helper itself harms the original agent. We therefore study an identification problem rather than a skill-learning problem: when does inserting a reusable temporal procedure actually repair a recurring failure? For each frozen endpoint, we compare a targeted procedure (T), a complexity-matched target-blind helper (G), and the original no-skill agent (N); a procedure is a binding repair under this intervention only when T beats both G and N. The extra anchor changes an observed scientific verdict: on one DeepSeek cutoff cell, T=100\% and G=72\%, but N=100\%, so the apparent +28-point two-arm gain is not a repair. Across 35 frozen endpoints reconstructed from first-party BEA, NOAA, and EIA release histories, positive effects appear selectively in non-ceiling regimes. The strongest held-out cutoff result reaches 100\% under T versus 70\% under G and 47.5\% under N, with endpoint sign test $p=0.0156$ for T--N; byte-identical procedures also retain direction on compatible future releases, while cross-domain grounding is null and a second model is already at no-skill ceiling on the same EIA tasks. The evidence supports a conditional diagnostic of reusable-procedure bottlenecks. It does not establish universal temporal-skill gains, behavioral neutrality of the executed generic helper, or superiority of a callable skill container over an equivalent retrieval-side implementation.
\end{abstract}
